\section{Conclusion}
PharmaScan combines a learned pill embedding, enrolment prototypes and imprint evidence to verify medicines with explicit rejection and retake decisions. Adaptation of a hybrid facial expression encoder enables enrolment of previously unseen medicines from three photographs. In the strict C3PI evaluation, the gallery using embeddings alone achieves 76.0\% correct naming and 0.04\% accepted misidentification among enrolled queries while rejecting 94.8\% of foreign queries, and imprint re-ranking increases catalogue top-1 accuracy from 71.6\% to 80.2\%. Together, the experiments show that useful embedding separation depends on the training recipe and that improved imprint recovery does not necessarily reduce accepted errors.

Acquisition transfer remains the main limitation because representation learning relies on C3PI consumer images and catalogue coverage is restricted to US products. On MEDISEG, rejection of 70.4\% of enrolled queries limits useful recognition despite rare misidentification among enrolled queries. Small pills and faint or debossed imprints further restrict the available evidence, with an imprint token recovered in about 47\% of validation photographs. The remaining ePillID performance gap and the absence of evaluation on real phone photographs of Singapore medicines and dispensing labels also delimit the current findings.

Future work centres on local enrolment with calibration for each patient or device, evaluating accepted errors alongside rejection and retakes. Controlled branch and fusion comparisons and the ePillID multi-head recipe extend the representation analysis, while adaptation of the imprint reader is assessed through text recovery and downstream verification errors. Prospective studies extend this evaluation to usability and medication errors.
